\documentclass[11pt]{article}
\usepackage{acl}

\usepackage{times}
\usepackage{latexsym}
\usepackage[T1]{fontenc}
\usepackage[utf8]{inputenc}
\usepackage{microtype}
\usepackage{inconsolata}
\usepackage{url}
\usepackage{booktabs}
\usepackage{amsmath}
\usepackage{amssymb}
\usepackage{enumitem}
\usepackage{graphicx}
\Urlmuskip=0mu plus 1mu
\emergencystretch=3em

\title{Predicting Market Outcomes from Earnings Call Transcripts:\\
       Fine-tuning, Retrieval, and Prompting}

\author{
Surya Suresh \\
\And
Nikhil Kumar \\
\And
Dheeraj Gosula \\
}

\begin{document}
\maketitle

\begin{abstract}
We study whether a modern open-weight large language model can predict
financial outcomes from S\&P~500 earnings call transcripts, and which
adaptation techniques actually help. We test the model on two binary
classification tasks: stock direction (whether the price moves up or
down over the three trading days after the call) and earnings surprise
(whether reported EPS beats analyst consensus). For each task we
evaluate three techniques: LoRA fine-tuning, retrieval-augmented
generation (RAG), and few-shot prompting. Using Qwen2.5-7B-Instruct as
a fixed base model, we run every combination of these techniques on
both tasks and report results across 16 conditions. LoRA fine-tuning
produces a substantial improvement on the surprise task (AUC 0.637 to
0.812, 95\% CI excluding random) but no measurable improvement on
stock direction (AUC remains near 0.50). Few-shot prompting and
retrieval show smaller, mixed effects: they tend to help the
unadapted base model but slightly hurt the fine-tuned one. The
asymmetry between the two tasks suggests that LLMs fine-tuned on
earnings calls can identify outcomes already discussed in the
transcript, but that predicting stock movement from text alone
remains close to chance.
\end{abstract}

\section{Introduction}

Earnings calls are quarterly events in which a public company's
management reports financial results and answers questions from
sell-side analysts. The accompanying transcripts are dense, structured
documents that contain explicit financial figures, forward-looking
guidance, and unscripted exchanges in the question-and-answer (Q\&A)
section. They are widely read by investors, but extracting predictive
signal from them is hard: short-horizon stock returns are noisy, and
the most informative content is mixed with boilerplate disclosures and
scripted pleasantries.

Most prior work on earnings call NLP relies on sentiment scoring or
single-task classifiers~\citep{Loughran2011,Araci2019,Qin2019}.
More recent work has applied larger transformer
architectures~\citep{Koval2023,Ma2020}, but no study has
systematically compared the three techniques most commonly used to
adapt large language models to a target domain---supervised
fine-tuning, retrieval-augmented generation
(RAG)~\citep{Lewis2020}, and few-shot
prompting~\citep{Brown2020,Wei2022}---on the same
transcripts and the same prediction tasks.

We close this gap with a $2^3$ factorial experiment. For each of
two tasks (stock direction and earnings surprise) we toggle
fine-tuning, RAG, and few-shot prompting independently and run all
eight resulting combinations. We use Qwen2.5-7B-Instruct~\citep{Qwen25}
as a fixed base model and adapt it efficiently with
QLoRA~\citep{Dettmers2023,Hu2022}, allowing the entire pipeline to run
on a single 16~GB V100 GPU.

Our contributions are:
\begin{itemize}[leftmargin=*,topsep=2pt,itemsep=2pt]
\item A controlled comparison of fine-tuning, retrieval, and prompting
for transcript-based outcome prediction, with bootstrap confidence
intervals on every reported metric.
\item Evidence that LoRA fine-tuning yields a large, statistically
meaningful gain on EPS surprise (\,$+0.18$ AUC over zero-shot base) while
leaving stock direction at chance.
\item A head-and-tail truncation strategy that fits within the V100
context budget while preserving both opening remarks and the Q\&A
section, two parts of the call that prior work flags as
high-signal~\citep{Qin2019,Cohen2020}.
\item Analysis of the asymmetry between the two tasks: surprise is
largely answerable from in-document linguistic and numerical cues,
whereas three-day stock direction depends on factors outside the
transcript.
\end{itemize}

\section{Related Work}

\paragraph{Sentiment-based approaches.}
The earliest text-based methods for financial prediction use sentiment
dictionaries such as Loughran-McDonald~\citep{Loughran2011} or
domain-specific classifiers like FinBERT~\citep{Araci2019,Yang2020}.
These models compress an entire transcript into a small number of
sentiment scores, which are fed to a downstream classifier. The
compression discards structural information, and as we confirm in
Section~\ref{sec:results}, reduces stock direction prediction to
roughly chance accuracy.

\paragraph{Earnings call deep learning.}
\citet{Qin2019} use multimodal text-and-audio features to predict
post-call return volatility. \citet{Ma2020} combine GloVe embeddings
with industry-sector signals to predict next-day stock movement and
report 48.5--56.8\% accuracy across sectors. \citet{Koval2023} apply
hierarchical FinBERT to earnings surprise classification and report
76.56\% accuracy on a balanced test set. \citet{Cohen2020} show that
the unscripted Q\&A portion of calls carries more predictive signal
than the scripted prepared remarks. Each of these works tests a single
architecture on a single task; none compare across modern
adaptation techniques.

\paragraph{Domain LLMs and parameter-efficient adaptation.}
BloombergGPT~\citep{Wu2023} pretrains a 50B-parameter model on
Bloomberg's proprietary financial corpus. We instead use a smaller,
openly available model (Qwen2.5-7B) and adapt it with
LoRA~\citep{Hu2022} and QLoRA~\citep{Dettmers2023}, which combines
4-bit weight quantization with low-rank trainable adapters and fits
fine-tuning on a single 16~GB GPU. Generic adaptation
techniques---retrieval~\citep{Lewis2020} and chain-of-thought
prompting~\citep{Wei2022}---are well studied in open-domain NLP but, to
our knowledge, untested side-by-side on earnings call outcome
prediction.

\section{Data}
\label{sec:data}

We use the publicly available S\&P~500 earnings transcripts corpus
\texttt{Bose345/\allowbreak sp500\_earnings\_\allowbreak transcripts}
on HuggingFace, which contains tens of thousands of calls. From this
corpus we sample the most recent 5{,}000 transcripts (collected
May~2024 -- May~2025) and attempt to fetch labels for each via the
\texttt{yfinance} Python library. Coverage differs by label type. For
stock direction, we fetch the closing price on the last trading day
before the call and on the third trading day after; the label is~1
(up) if the post-call price exceeds the pre-call price, 0~(down)
otherwise. 1{,}965 calls yield valid direction labels. For earnings
surprise, we use yfinance's \texttt{get\_earnings\_dates()} method to
retrieve reported EPS and analyst consensus, requiring a match within
10 days of the call; the label is~1 (beat) if reported $\geq$
consensus, 0~(miss) otherwise. 827 calls yield valid surprise labels.

The sharply lower coverage for surprise reflects yfinance's limited
EPS data for delisted, recently-IPO'd, or non-standard tickers
(class-B dual-share entries, recent corporate actions, etc.), and is
the single largest external constraint on our experimental scale.
Table~\ref{tab:dataset} summarizes the resulting datasets. Direction
is nearly class-balanced; surprise is heavily skewed toward
beat---reflecting the well-known tendency of S\&P 500 firms to beat
consensus more often than they miss. We mitigate this during training
(Section~\ref{sec:finetune}) but preserve the natural distribution at
test time so that reported metrics reflect deployment conditions.

\begin{table}[t]
\centering
\small
\begin{tabular}{lrrrr}
\toprule
\textbf{Task} & \textbf{Total} & \textbf{Train} & \textbf{Test} & \textbf{\% pos.} \\
\midrule
Direction & 1{,}965 & 1{,}768 & 197 & 51.1 \\
Surprise  &     827 &     744 &  83 & 80.9 \\
\bottomrule
\end{tabular}
\caption{Dataset sizes after fetching labels. Train/test split is
chronological 90/10 (no shuffle). \emph{\% pos.} is the fraction
labeled \emph{up} (direction) or \emph{beat} (surprise).}
\label{tab:dataset}
\end{table}

For each task we sort calls by date and use the earliest 90\% for
training and the latest 10\% for testing. A random split would leak
information from later calls into earlier predictions and inflate test
performance, particularly for direction where the macro market regime
itself carries non-trivial signal. Chronological splitting prevents
this lookahead bias and matches the deployment setting in which a
model trained on past calls must predict on new ones.

Transcripts are long: the median length under the Qwen tokenizer is
approximately 13{,}000 tokens, far exceeding our 2{,}048-token budget
per training example. Naive front-truncation (keep first 2{,}048
tokens) preserves the prepared opening remarks but discards the entire
Q\&A section. The Q\&A is unscripted, contains analyst pressure on
specific concerns, and is widely cited as a high-signal segment of the
call~\citep{Qin2019,Cohen2020}. We instead use a head-and-tail
truncation: keep the first 60\% (1{,}229 tokens) and the last 40\%
(819~tokens) of the budget, dropping only the middle. This preserves
both the management narrative and the Q\&A. The same truncation is
used at training and at evaluation so that the prompt distribution
seen by the model is consistent across the two phases.

\section{Methodology}

\subsection{Task formulation}
\label{sec:formulation}

Each prompt is constructed as
\begin{verbatim}
<system instruction>

Transcript:
<truncated transcript>

Answer:
\end{verbatim}
where \texttt{<system instruction>} is a one-paragraph description of
the task with the allowed answer tokens (\emph{up}/\emph{down} or
\emph{beat}/\emph{miss}). The model is asked to emit one token after
\texttt{Answer:}.

At inference we run a single forward pass and read the logits at the
final position. We extract the two target token IDs (e.g.,
\texttt{\char32 up} and \texttt{\char32 down}) and apply softmax over
just that pair to obtain a probability of the positive class. This
avoids generation entirely: we never sample from the model and never
parse free-form text. The same scoring procedure is used for the
zero-shot base model and for all augmented variants, so all conditions
are directly comparable.

\subsection{Fine-tuning with QLoRA}
\label{sec:finetune}

We fine-tune Qwen2.5-7B-Instruct~\citep{Qwen25} using
QLoRA~\citep{Dettmers2023}, which combines 4-bit weight quantization
with low-rank adapter training~\citep{Hu2022}. The base model is
loaded with 4-bit NF4 quantization and double quantization, reducing
weight memory from $\sim$14~GB (fp16) to $\sim$4~GB---enough headroom
to fit the 7B model on a 16~GB V100 along with activations and
gradients. We attach LoRA adapters to the four attention projection
matrices in every transformer block (\texttt{q\_proj},
\texttt{k\_proj}, \texttt{v\_proj}, \texttt{o\_proj}) at rank $r=16$,
scaling $\alpha=32$, and dropout 0.05. Feed-forward and embedding
layers are entirely frozen, yielding approximately 10~M trainable
parameters out of 7.6~B total (0.13\%). Targeting attention
projections is the standard recipe in the LoRA literature: attention
controls how tokens combine information across positions, which is
the locus of any change we want a domain-adapted model to learn
(e.g., attending more strongly to Q\&A guidance language).

We train each task's adapter independently for 3~epochs with the
AdamW optimizer in 8-bit paged form. Learning rate is 2e-4 with
cosine decay and 3\% linear warmup, per-device batch size is~1 with
gradient accumulation over 8~steps for an effective batch size of~8,
and mixed-precision is fp16 (V100 hardware does not support bf16).
We enable gradient checkpointing to fit activations within memory and
apply a label mask so that the cross-entropy loss is computed only on
the answer token. For surprise (81\% beat / 19\% miss) the natural
class skew causes the model to collapse onto the majority class, so
we oversample the minority class within the training split only by
sampling miss examples with replacement until both classes have equal
counts; the validation and test splits are left at their natural
distribution. We evaluate on the chronological validation set after
each epoch and retain the checkpoint with lowest validation loss,
which guards against the onset of overfitting we observe after
epoch~1 on the smaller surprise dataset.

Each fine-tuned model is saved as a $\sim$40~MB LoRA adapter
alongside the saved tokenizer, decoupled from the unchanged base
weights. Adapters are applied at inference by loading the base model
with the same 4-bit quantization and attaching the adapter via the
PEFT library; the same procedure is used for the RAG and few-shot
variants that build on top of the fine-tuned model.

\subsection{Few-shot prompting}
\label{sec:prompt}

\textit{[Section to be completed by teammate.]} The few-shot variant
prepends $k$ labeled examples drawn chronologically from the training
split before the test transcript. Per the team's design, $k=4$ shots
are used, balanced across labels (2~positive, 2~negative), with
150~tokens per shot and a 3072-token total prompt budget.
Chain-of-thought reasoning was tested but did not yield reliable
improvements in our setting and is excluded from the final pipeline,
likely because the additional reasoning step disturbs the
single-token answer format the adapter was trained on. Extractive
sentence selection using a 24-keyword financial dictionary is used to
densify each shot's transcript content. Final scoring uses the same
softmax-over-target-pair procedure described in
Section~\ref{sec:formulation}.

\subsection{Retrieval-augmented generation}
\label{sec:rag}

\textit{[Section to be completed by teammate.]} The retrieval corpus
contains $\sim$450 documents, one per training transcript, each
consisting of the company ticker, call date, and first 2{,}000
characters of the transcript. Retrieval is sparse: the system first
attempts a same-company prior-quarter lookup (most recent transcripts
from the same ticker with date strictly preceding the query call),
then falls back to BM25 keyword ranking over the full training corpus
to fill remaining slots. Retrieved passages are prepended to the
prompt within a $\sim$400-token budget. The corpus is restricted to
the training split to prevent any test leakage.

\section{Experimental Setup}
\label{sec:setup}

All training and evaluation runs use a single NVIDIA V100~PCIE 16~GB
GPU on the OSC Pitzer cluster, with the Python environment pinned to
PyTorch~2.3.1+cu118 (the latest version supporting the V100's sm\_70
architecture), CUDA 11.8, transformers~4.44.2, peft~0.13.2, and
bitsandbytes~0.43.3.

For each task we evaluate eight conditions formed by toggling
fine-tuning, RAG, and few-shot prompting: zero-shot base, fine-tuned
only, RAG only, prompting only, and the four pairwise and triple
combinations. The base-only and prompting-only variants share a single
set of base-model weights; the fine-tuned variants use a task-specific
adapter on top of the same base model.

For each condition we report three metrics. The plain
prediction-vs-label match at threshold 0.5 (accuracy) is reported but
not relied on for claims, since on imbalanced tasks it is dominated
by the class prior. Balanced accuracy is the mean of per-class
recall, $(\mathrm{recall}_+ + \mathrm{recall}_-)/2$, which is robust
to class imbalance. AUC-ROC with a 95\% bootstrap confidence interval
(1{,}000 resamples) is our primary metric: it is threshold-free,
robust to imbalance, and the CI quantifies how much of the gap to
chance is statistically real given the test-set size. We treat a
result as above chance only if its 95\% AUC CI excludes 0.50, and
treat one model as better than another only when their CIs are
clearly separated; otherwise we report the gap with the caveat that
it is within sampling noise.

\section{Results}
\label{sec:results}

\subsection{Fine-tuning vs.\ zero-shot base}

Table~\ref{tab:ft-results} presents the headline fine-tuning results.

\begin{table}[t]
\centering
\small
\setlength{\tabcolsep}{4pt}
\begin{tabular}{llcccc}
\toprule
\textbf{Task} & \textbf{Model} & \textbf{Acc} & \textbf{B-Acc} & \textbf{AUC} & \textbf{95\% CI} \\
\midrule
Direction & Base       & 0.569 & 0.502 & 0.493 & [0.41, 0.57] \\
Direction & Fine-tuned & 0.594 & 0.500 & 0.499 & [0.42, 0.58] \\
\midrule
Surprise  & Base       & 0.759 & 0.566 & 0.637 & [0.49, 0.78] \\
Surprise  & \textbf{FT} & \textbf{0.807} & \textbf{0.708} & \textbf{0.812} & \textbf{[0.69, 0.92]} \\
\bottomrule
\end{tabular}
\caption{Fine-tuning vs.\ zero-shot base. Direction $n=197$,
surprise $n=83$. \emph{B-Acc} is balanced accuracy. AUC CIs are
bootstrap, 1{,}000 resamples.}
\label{tab:ft-results}
\end{table}

On the surprise task, fine-tuning improves AUC by 0.18
(0.637\,$\rightarrow$\,0.812) and balanced accuracy by 0.14
(0.566\,$\rightarrow$\,0.708). The fine-tuned AUC's 95\% CI
\mbox{[0.69,\,0.92]} cleanly excludes random chance (0.50). The two
CIs overlap by a small margin in the \mbox{[0.69,\,0.78]} window, so
the base-vs-fine-tuned gap is best described as substantial and
likely real, but at the edge of statistical separation given $n=83$.

On the direction task, both base and fine-tuned models hover near
AUC 0.50 with overlapping CIs. The fine-tuned model's plain accuracy
(0.594) matches the majority-class baseline
(117/197\,$\approx$\,0.594) almost exactly, indicating the model has
effectively learned to predict ``up'' on every example. Balanced
accuracy is exactly 0.500, confirming this collapse. Despite
minority-class oversampling during training, the model could not
extract any directional signal beyond the prior.

These results are broadly consistent with prior work.
\citet{Ma2020} report 48.5--56.8\% accuracy across sectors for
next-day stock movement using GloVe and industry features---similar
to our chance-level direction result, despite their use of explicit
sector signals. \citet{Koval2023} report 76.56\% accuracy on a
balanced earnings surprise test set with hierarchical FinBERT; our
fine-tuned model reaches 80.7\% accuracy on a 77/19 imbalanced test
and AUC 0.812 on the same test, suggesting that decoder-LLM
fine-tuning is at least competitive with task-specialized encoder
architectures on this task. We do not claim a head-to-head
improvement since the test sets and label-fetching protocols differ.

\subsection{Few-shot prompting}

\textit{[Analysis to be completed by teammate.]} The 4-shot prompting
results from our presentation slides are summarized in
Table~\ref{tab:prompt-stub} for completeness; the discussion of
\emph{why} the effect of prompting is asymmetric across base and
fine-tuned models, and across the two tasks, is left to the teammate
who designed the prompting pipeline.

\begin{table}[h]
\centering
\small
\setlength{\tabcolsep}{4pt}
\begin{tabular}{lcc}
\toprule
\textbf{Condition} & \textbf{Zero-shot AUC} & \textbf{4-shot AUC} \\
\midrule
Base direction & 0.493 & 0.520 \\
FT direction   & 0.499 & 0.536 \\
Base surprise  & 0.637 & 0.720 \\
FT surprise    & 0.812 & 0.748 \\
\bottomrule
\end{tabular}
\caption{Few-shot prompting AUC, copied from team presentation.
Confidence intervals are reported in the original slide deck.}
\label{tab:prompt-stub}
\end{table}

\subsection{Retrieval-augmented generation}

\textit{[Analysis to be completed by teammate.]}
Table~\ref{tab:rag-stub} summarizes the RAG results from the team
presentation; analysis of why RAG helps the base model on surprise but
slightly hurts the fine-tuned model is left to the teammate.

\begin{table}[h]
\centering
\small
\setlength{\tabcolsep}{4pt}
\begin{tabular}{lcccc}
\toprule
\textbf{Cond.} & \textbf{ZS} & \textbf{4-shot} & \textbf{RAG} & \textbf{RAG+4} \\
\midrule
Base direction & 0.493 & 0.520 & 0.492 & 0.532 \\
FT direction   & 0.499 & 0.536 & 0.494 & 0.538 \\
Base surprise  & 0.637 & 0.720 & 0.680 & 0.676 \\
FT surprise    & 0.812 & 0.748 & 0.752 & 0.722 \\
\bottomrule
\end{tabular}
\caption{Combined results across all four conditions. ZS = zero-shot.
Numbers are AUC; full CIs are reported in the original slide deck.}
\label{tab:rag-stub}
\end{table}

\section{Discussion}

\subsection{Why does fine-tuning help one task but not the other?}

The asymmetry between surprise and direction in
Table~\ref{tab:ft-results} is not a quirk of model capacity or training
procedure---it reflects a fundamental difference in where the
predictive signal lives.

For surprise, the relevant information is largely contained in the
transcript itself. The earnings press release that announces actual
EPS is published roughly 30~minutes before the call begins; by the
time the CFO reads opening remarks, the beat or miss is already
public, and management almost always references the figures
explicitly (e.g., ``We delivered adjusted EPS of \$1.85, ahead of
consensus of \$1.77''). Even when our truncation drops the precise
numbers, the framing, tone, guidance language, and Q\&A dynamics
provide strong proxies. Fine-tuning teaches the model which proxies to
attend to.

For direction, the label is determined by post-call market behavior
over three trading days---a function of how thousands of investors
interpret the call and of intervening macro news, rate moves, sector
rotations, and quarterly fund flows. None of that is in the
transcript. The model is being asked to forecast something whose
causal driver is largely outside the input. We should expect text-only
models to fail at this, and they do.

This is consistent with the broader literature on post-earnings
announcement drift: even with audio features and structured market
data, prior work reports only modest gains over chance on
short-horizon direction~\citep{Qin2019,Sawhney2020}, and our
text-only setup is strictly less informed.

\subsection{The role of head-and-tail truncation}

Several earnings call NLP studies note that the Q\&A section carries
more candid signal than the prepared remarks~\citep{Qin2019,Cohen2020},
since management's scripted opening tends to be uniformly positive
regardless of actual performance. Our head-and-tail strategy
operationalizes this: at the same 2{,}048-token budget, we retain both
segments rather than discarding the Q\&A entirely. We did not run a
controlled ablation against front-truncation, but our 0.812 AUC is at
the high end of comparable BERT-family results on similar transcript
spans, suggesting the truncation choice is consequential. A natural
follow-up experiment is to repeat the surprise fine-tune with strict
front-truncation under matched conditions and quantify the gap.

\subsection{Class imbalance and metric choice}

The surprise task's 81/19 class skew is a recurring trap. A model that
predicts ``beat'' on every example achieves $\sim$77\% accuracy with
zero predictive skill, which is dangerously close to our base model's
0.759 accuracy. We rely on AUC and balanced accuracy precisely because
they cannot be inflated by majority-class collapse. The direction task
illustrates the same point in reverse: the fine-tuned model's plain
accuracy (0.594) looks like a small improvement over base (0.569), but
its balanced accuracy (0.500) reveals that all gain comes from
learning the prior more confidently, not from any new discrimination
ability.

\subsection{Effect of method combination}

Reading across Tables~\ref{tab:ft-results}--\ref{tab:rag-stub}, two
patterns emerge. First, on the surprise task, every additional
augmentation on top of fine-tuning slightly degrades AUC
(0.812\,$\rightarrow$\,0.748 with 4-shot;
0.812\,$\rightarrow$\,0.752 with RAG;
0.812\,$\rightarrow$\,0.722 with both). The fine-tuned adapter has
already learned what to look for from the in-document signal, and
adding noisier context appears to crowd out the signal-rich Q\&A
within the fixed context window. Second, on the surprise task without
fine-tuning, augmentations help (4-shot: 0.637\,$\rightarrow$\,0.720;
RAG: 0.637\,$\rightarrow$\,0.680). For an unadapted base model more
context is better; for an adapted one, the context budget is better
spent on the test transcript itself. This is consistent with prior
findings that retrieval and in-context learning provide diminishing
returns once a task has been fine-tuned.

For direction, no condition meaningfully exceeds chance, suggesting
the bottleneck is signal availability in the input, not modeling
choices.

\section{Limitations}

\paragraph{Temporal coverage.}
Our entire dataset spans roughly one year (May~2024 -- May~2025). This
covers a single market regime and only a few earnings cycles.
Performance under different macro conditions (e.g., a correction, a
high-volatility episode, or rate-cycle inflection) is unknown.
Expanding to multi-year data would improve external validity and would
likely tighten the surprise CIs by enlarging the test set.

\paragraph{Test set size.}
The surprise test set ($n=83$, of which 19 are \emph{miss}) is small
enough that AUC confidence intervals are wide ($\pm 0.11$). This is
driven by yfinance's limited EPS coverage---the analyst-consensus
endpoint succeeds on only $\sim$40\% of S\&P 500 tickers in our
sampled window. A different EPS data source (FactSet, Refinitiv,
yahooquery) would likely double the labeled dataset.

\paragraph{Single base model.}
We did not run an ablation across base LLMs. Llama-3-8B,
Gemma-2-9B, and finance-pretrained models like FinMA may behave
differently. The qualitative conclusion---that fine-tuning helps
surprise but not direction---would likely hold (the
signal-availability story is model-agnostic), but specific AUC numbers
would shift.

\paragraph{Truncation budget.}
Our 2{,}048-token cap is a hardware constraint of the V100. Extending
to 4096 or 8192 tokens would let us include more middle prepared
remarks and more Q\&A content; this is a natural follow-up on an A100
or larger GPU.

\paragraph{No tradeable edge.}
Surprise prediction from transcripts is a useful capability test, but
not a tradeable signal. By the time the call begins, the press release
with actual EPS numbers has been public for $\sim$30~minutes; the
``prediction'' is in fact a reading-comprehension task. We frame our
results accordingly: the surprise gain is evidence of financial
language understanding, not of forecasting capability.

\section{Conclusion}

We tested whether a modern open-weight LLM, adapted via fine-tuning,
retrieval, or prompting, can predict financial outcomes from earnings
call transcripts. LoRA fine-tuning produces a clear and statistically
meaningful gain on EPS surprise (AUC 0.812 vs.\ 0.637 base) but does
not move stock direction prediction off chance. The asymmetry is
interpretable: surprise is largely answerable from in-document
evidence, while three-day stock direction depends on factors absent
from the transcript. Combining augmentations on top of a fine-tuned
adapter provides no further gain on surprise and often slightly
hurts. Our results support the use of fine-tuned LLMs for
transcript-grounded financial classification tasks, and caution
against expecting them to forecast short-horizon market reactions
from text alone.

\section*{Author Contributions}

\paragraph{Surya Suresh.}
% TODO: fill in -- e.g., few-shot prompting design and analysis,
% extractive sentence selection, prompting results.

\paragraph{Nikhil Kumar.}
% TODO: fill in -- e.g., data pipeline, boilerplate filtering,
% QLoRA fine-tuning, evaluation metrics, infrastructure on
% OSC Pitzer/Ascend/Cardinal.

\paragraph{Dheeraj Gosula.}
% TODO: fill in -- e.g., retrieval-augmented generation pipeline,
% sparse retrieval corpus, BM25 fallback, RAG results.

\section*{Acknowledgements}

AI assistants (Claude and ChatGPT) were used to help organize and
format the LaTeX code for this report and as a mild writing
assistant.

\section*{References}
\begingroup
\renewcommand{\section}[2]{}
\begin{thebibliography}{}

\bibitem[Araci(2019)]{Araci2019}
Dogu Araci. 2019.
FinBERT: Financial Sentiment Analysis with Pre-trained Language Models.
\textit{arXiv preprint arXiv:1908.10063}.

\bibitem[Brown et al.(2020)]{Brown2020}
Tom Brown, Benjamin Mann, Nick Ryder, et al. 2020.
Language Models are Few-Shot Learners.
In \textit{NeurIPS}.

\bibitem[Cohen et al.(2020)]{Cohen2020}
Lauren Cohen, Christopher Malloy, and Quoc Nguyen. 2020.
Lazy Prices.
\textit{Journal of Finance}, 75(3):1371--1415.

\bibitem[Dettmers et al.(2023)]{Dettmers2023}
Tim Dettmers, Artidoro Pagnoni, Ari Holtzman, and Luke Zettlemoyer. 2023.
QLoRA: Efficient Finetuning of Quantized LLMs.
In \textit{NeurIPS}.

\bibitem[Hu et al.(2022)]{Hu2022}
Edward Hu, Yelong Shen, Phillip Wallis, et al. 2022.
LoRA: Low-Rank Adaptation of Large Language Models.
In \textit{ICLR}.

\bibitem[Koval et al.(2023)]{Koval2023}
Ross Koval, Nicholas Andrews, and Xifeng Yan. 2023.
Forecasting Earnings Surprises from Conference Call Transcripts.
In \textit{Findings of ACL}.

\bibitem[Lewis et al.(2020)]{Lewis2020}
Patrick Lewis, Ethan Perez, Aleksandra Piktus, et al. 2020.
Retrieval-Augmented Generation for Knowledge-Intensive NLP Tasks.
In \textit{NeurIPS}.

\bibitem[Loughran and McDonald(2011)]{Loughran2011}
Tim Loughran and Bill McDonald. 2011.
When Is a Liability Not a Liability? Textual Analysis, Dictionaries, and 10-Ks.
\textit{Journal of Finance}, 66(1):35--65.

\bibitem[Ma et al.(2020)]{Ma2020}
Yu Ma, Rui Mao, Qika Lin, et al. 2020.
Multi-Source Aggregated Classification for Stock Price Movement Prediction.
\textit{Information Fusion}.

\bibitem[Qin and Yang(2019)]{Qin2019}
Yao Qin and Yi Yang. 2019.
What You Say and How You Say It Matters: Predicting Stock Volatility Using Verbal and Vocal Cues.
In \textit{ACL}, pages 390--401.

\bibitem[Qwen Team(2024)]{Qwen25}
Qwen Team. 2024.
Qwen2.5 Technical Report.
\textit{arXiv preprint arXiv:2412.15115}.

\bibitem[Sawhney et al.(2020)]{Sawhney2020}
Ramit Sawhney, Piyush Khanna, Arshiya Aggarwal, Taru Jain, Puneet Mathur, and Rajiv Ratn Shah. 2020.
VolTAGE: Volatility Forecasting via Text Audio Fusion with Graph Convolution Networks for Earnings Calls.
In \textit{EMNLP}.

\bibitem[Wei et al.(2022)]{Wei2022}
Jason Wei, Xuezhi Wang, Dale Schuurmans, et al. 2022.
Chain-of-Thought Prompting Elicits Reasoning in Large Language Models.
In \textit{NeurIPS}.

\bibitem[Wu et al.(2023)]{Wu2023}
Shijie Wu, Ozan Irsoy, Steven Lu, et al. 2023.
BloombergGPT: A Large Language Model for Finance.
\textit{arXiv preprint arXiv:2303.17564}.

\bibitem[Yang et al.(2020)]{Yang2020}
Yi Yang, Mark Christopher Siy Uy, and Allen Huang. 2020.
FinBERT: A Pretrained Language Model for Financial Communications.
\textit{arXiv preprint arXiv:2006.08097}.

\end{thebibliography}
\endgroup

\end{document}
